\documentclass{article}
\begin{document}
\section*{Positive chain categories in the presented span powers}
Fix a possibly nondiscrete base, a strict category-valued presheaf of grades and the existing split endpoint data. For each positive length the left-associated Grothendieck span presentation has an apex functor to independent chains. It is an equivalence of categories. Both composites with the source and destination functors are equal to the corresponding native leg, not merely equal on objects. No general grade arrow or endpoint lift is required to be invertible.
\par The last-grade construction casts the final arrow across both endpoint equalities and the prefix arrow comparison, controlling its outgoing map by the successor formula. The bijection restricts to squares over each prescribed incoming base arrow. Identity and pasting use the headwise rebased-tail laws for the independent final-grade decomposition. For native pasting, equality of the prefix alone is insufficient: the proof compares the reindexed final grade composites and their trailing transports. The category of chains uses the pre-existing independent square identity and pasting; only the proofs of the category axioms appeal to the proved bijection and preservation laws.
\par Positive index zero means length one. At each successor the native presentation is equivalent to the strict matching pullback; it is not asserted to be literally that pullback. Not proved: a strict inverse commuting with both legs, reassociation between all bracketings, the empty span, extension of the identified equipment span against every payload, the coproduct and monoidality step, flattening, strict algebras, or closure of issue 105. The last-grade argument is not a reversal of a noncommutative word.
\end{document}
